\subsection{Derivations and liftings of homomorphisms}

Let \(k\) be a ring, C a \(k\)-algebra, and N an ideal of C with square zero. Denote by \(\pi : \mathrm{C} \to \mathrm{C/N}\) the canonical homomorphism; since \(\mathrm{N}^2 = \{0\}\), the C-module structure of N comes from a C/N-module structure.

Let A be a \(k\)-algebra and \(\varphi : \mathrm{A} \to \mathrm{C/N}\) a homomorphism of \(k\)-algebras. Equip N with the A-module structure deduced from \(\varphi\). We call a \(lifting\) of \(\varphi\) (to C) every homomorphism of \(k\)-algebras \(\tilde{\varphi} : \mathrm{A} \to \mathrm{C}\) such that \(\pi \circ \tilde{\varphi} = \varphi\). Let \(\tilde{\varphi}\) be such a lifting, and let \(\delta\) be a map from A into N; denote by \(\delta + \tilde{\varphi}\) the map \(x \mapsto \delta(x) + \tilde{\varphi}(x)\) from \(A\) into C.

PROPOSITION 1.--- If φ admits a lifting, the map (δ, φ) ↦ δ + φ defines a simply transitive action of the group of k-derivations of A into N on the set of liftings of φ.

Let \(\tilde{\varphi}_{0}: A \to C\) be a lifting of \(\varphi\) . The map \(\delta \mapsto \delta + \tilde{\varphi}_{0}\) induces a bijection from the set of maps from \(A\) into \(N\) onto the set of maps \(\tilde{\varphi}: A \to C\) such that \(\pi \circ \tilde{\varphi} = \varphi\) . Fix \(\delta\), and set \(\tilde{\varphi} = \delta + \tilde{\varphi}_{0}\) For \(\tilde{\varphi}\) to be a homomorphism of \(k\) -algebras, it is necessary and sufficient that \(\delta\) be a \(k\)-derivation from \(A\) into \(N\) : indeed, for \(x, y\) in \(A\) and \(\lambda\) in \(k\), we have the relations

\[
\begin{array}{c} \tilde {\varphi} (x + y) - \tilde {\varphi} (x) - \tilde {\varphi} (y) = \delta (x + y) - \delta (x) - \delta (y) \\ \tilde {\varphi} (\lambda x) - \lambda \tilde {\varphi} (x) = \delta (\lambda x) - \lambda \delta (x) \\ \tilde {\varphi} (x y) - \tilde {\varphi} (x) \tilde {\varphi} (y) = \delta (x y) - \delta (x) \delta (y) - \delta (x) \tilde {\varphi} _ {0} (y) - \tilde {\varphi} _ {0} (x) \delta (y) \\ = \delta (x y) - \varphi (x) \delta (y) - \varphi (y) \delta (x), \end{array}
\]

the last equality follows from the fact that N has square zero. The proposition follows.

Example.--- Let B be a k-algebra, N a B-module. Equip the k-module B⊕N with the k-algebra structure defined by \((b,x)(b',x') = (bb',bx' + b'x)\) (cf. A, III, p. 127), so that N is an ideal of square zero in B ⊕ N. Let \(\varphi : A \to B\) be a homomorphism of k-algebras. Then the liftings of \(\varphi\) to B ⊕ N are the maps \(x \mapsto (\varphi(x), \delta(x))\) , where \(\delta\) ranges over the set of k-derivations of A into N (cf. loc. cit., Proposition 12).

Let \(\Omega_{k}(A)\) be the module of \(k\)-differentials of the ring A, and let \(d: \mathrm{A} \longrightarrow \Omega_{k}(\mathrm{A})\) be the universal \(k\)-derivation (A, III, pp. 133–134); recall (loc. cit.) that for every A-module M, the map \(v \mapsto v \circ d\) is an A-linear isomorphism of \(\mathrm{Hom}_{\mathrm{A}}(\Omega_{k}(A), \mathrm{M})\) onto the A-module of \(k\)-derivations of A into M.

Let J be an ideal of A. According to A, III, p. 137, we have an exact sequence of A/J-linear maps

\[
\mathrm{J} / \mathrm{J} ^ {2} \stackrel {\bar {d}} {\longrightarrow} (\mathrm{A} / \mathrm{J}) \otimes_ {\mathrm{A}} \Omega_ {k} (\mathrm{A}) \longrightarrow \Omega_ {k} (\mathrm{A} / \mathrm{J}) \rightarrow 0,
\]

where \(\bar{d}\) is the homomorphism deduced by passing to quotients from the restriction of \(d\) to J.

Let \(\rho : A \to A/J^{2}\) and \(\pi : A/J^{2} \to A/J\) the canonical surjections. Let \(v : (A/J) \otimes_{A} \Omega_{k}(A) \longrightarrow J/J^{2}\) a \(k\)-linear; associate to it a \(k\)-linear \(H_{v} : A \to A/J^{2}\) by setting \(H_{v}(x) = \rho(x) - v(1 \otimes dx)\) If \(v\) is a retraction of \(\bar{d}\), \(H_{v}\) vanishes on \(J\) , hence induces by passage to the quotient a \(k\)-linear \(h_{v} : A/J \to A/J^{2}\) . On the other hand, given a \(k\)-linear \(h : A/J \to A/J^{2}\) denote \(\psi_{h} : A/J \oplus J/J^{2} \longrightarrow A/J^{2}\) the map \((x,y) \mapsto h(x)+y\) .

PROPOSITION 2. - Equip the \(k\)-module \(A / J \oplus J / J^2\) with the structure of \(k\)-algebra defined in the example above. The maps \(v \mapsto h_v\) and \(h \mapsto \psi_h\) induce bijections between the following sets:

\begin{enumerate}
\def\labelenumi{(\roman{enumi})}
\item
  the set of A/J-linear retractions v of \(\bar{d}\) ;
\item
  the set of homomorphisms of k-algebras \(h: A/J \to A/J^{2}\) such that \(\pi \circ h = Id_{A/J}\) ;
\item
  the set of isomorphisms of k-algebras \(\psi : \mathrm{A} / \mathrm{J} \oplus \mathrm{J} / \mathrm{J}^2 \longrightarrow \mathrm{A} / \mathrm{J}^2\) such that \(\pi \circ \psi = \mathrm{pr}_1\) and \(\psi(0, z) = z\) for \(z \in \mathrm{J} / \mathrm{J}^2\) .
\end{enumerate}

Let us apply Proposition 1 with \(C = A / J^2\) and \(N = J / J^2\). Let \(\varphi : A \to A / J\) be the canonical surjection; the homomorphism \(\rho\) is a lifting of \(\varphi\) to \(A / J^2\). The A-module \(\text{Hom}_{A / J}((A / J) \otimes_A \Omega_k(A), J / J^2)\) is identified with \(\text{Hom}_A(\Omega_k(A), J / J^2)\) ; by Proposition 1, the map \(v \mapsto H_v\) is a bijection from this set onto the set of liftings of \(\varphi\) to \(A / J^2\) To show that \(x \in J\) , one has \(1 \otimes dx = \bar{d}(\rho(x))\) ; for \(H_v\) to vanish on \(J\) , it is necessary and sufficient that \(v \circ \bar{d}\) be the identity map of \(J / J^2\). This proves that the map \(v \mapsto h_v\) induces a bijection between the first two sets described in the statement.

The map \(h \mapsto \psi_h\) is a bijection from the set of \(k\)-linear maps from A/J to A/J \(^2\) onto the set of \(k\)-linear maps \(\psi : \text{A/J} \oplus \text{J/J}^2 \longrightarrow \text{A/J}^2\) such that \(\psi(0,z) = z\) for \(z \in \text{J/J}^2\); moreover, for \(\pi \circ \psi_h = \text{pr}_1\) it is necessary and sufficient that \(\pi \circ h = \text{Id}_{\text{A/J}}\), that is \(z \equiv h(\pi(z)) \pmod{\text{J/J}^2}\) for every \(z \in \text{A/J}^2\). Suppose these conditions are satisfied. For \(h\) to be a ring homomorphism, it is necessary and sufficient that the same be true of \(\psi_h\); moreover, the homomorphism \(\psi_h\) is bijective: the inverse map associates to an element \(z\) of A/J \(^2\) the pair \((\pi(z), z - h(\pi(z)))\). This proves that the map \(h \mapsto \psi_h\) induces a bijection between the last two sets described in the statement.

